\documentclass[11pt, a4paper]{article}
%\usepackage{geometry}
\usepackage[inner=1.5cm,outer=1.5cm,top=2.5cm,bottom=2.5cm]{geometry}
\pagestyle{empty}
\usepackage{graphicx}
\usepackage{fancyhdr, lastpage, bbding, pmboxdraw}
\usepackage[usenames,dvipsnames]{color}
\definecolor{darkblue}{rgb}{0,0,.6}
\definecolor{darkred}{rgb}{.7,0,0}
\definecolor{darkgreen}{rgb}{0,.6,0}
\definecolor{red}{rgb}{.98,0,0}
\usepackage[colorlinks,pagebackref,pdfusetitle,urlcolor=darkblue,citecolor=darkblue,linkcolor=darkred,bookmarksnumbered,plainpages=false]{hyperref}
\renewcommand{\thefootnote}{\fnsymbol{footnote}}

\pagestyle{fancyplain}
\fancyhf{}
\lhead{ \fancyplain{}{Advanced Topics in Network Science} }
%\chead{ \fancyplain{}{} }
\rhead{ \fancyplain{}{\today} }
%\rfoot{\fancyplain{}{page \thepage\ of \pageref{LastPage}}}
\fancyfoot[RO, LE] {page \thepage\ of \pageref{LastPage} }
\thispagestyle{plain}

%%%%%%%%%%%% LISTING %%%
\usepackage{listings}
\usepackage{caption}
\DeclareCaptionFont{white}{\color{white}}
\DeclareCaptionFormat{listing}{\colorbox{gray}{\parbox{\textwidth}{#1#2#3}}}
\captionsetup[lstlisting]{format=listing,labelfont=white,textfont=white}
\usepackage{verbatim} % used to display code
\usepackage{fancyvrb}
\usepackage{acronym}
\usepackage{amsthm}
% \VerbatimFootnotes overflows TeX's input stack on the Attendance item, whose
% three consecutive \footnote{\url{...}} are more than it can re-scan. No
% footnote here contains verbatim, so the package hook buys nothing; leaving it
% enabled costs the whole PDF.
%\VerbatimFootnotes



\definecolor{OliveGreen}{cmyk}{0.64,0,0.95,0.40}
\definecolor{CadetBlue}{cmyk}{0.62,0.57,0.23,0}
\definecolor{lightlightgray}{gray}{0.93}



\lstset{
%language=bash,                          % Code langugage
basicstyle=\ttfamily,                   % Code font, Examples: \footnotesize, \ttfamily
keywordstyle=\color{OliveGreen},        % Keywords font ('*' = uppercase)
commentstyle=\color{gray},              % Comments font
numbers=left,                           % Line nums position
numberstyle=\tiny,                      % Line-numbers fonts
stepnumber=1,                           % Step between two line-numbers
numbersep=5pt,                          % How far are line-numbers from code
backgroundcolor=\color{lightlightgray}, % Choose background color
frame=none,                             % A frame around the code
tabsize=2,                              % Default tab size
captionpos=t,                           % Caption-position = bottom
breaklines=true,                        % Automatic line breaking?
breakatwhitespace=false,                % Automatic breaks only at whitespace?
showspaces=false,                       % Dont make spaces visible
showtabs=false,                         % Dont make tabls visible
columns=flexible,                       % Column format
morekeywords={__global__, __device__},  % CUDA specific keywords
}

%%%%%%%%%%%%%%%%%%%%%%%%%%%%%%%%%%%%
\begin{document}
\begin{center}
{\Large \textsc{Advanced Topics in Network Science}}
\end{center}
\begin{center}
Fall 2026
\end{center}
%\date{September 26, 2014}

\begin{center}
\rule{6in}{0.4pt}
\begin{minipage}[t]{.75\textwidth}
\begin{tabular}{llcccll}
\textbf{Instructor:} & Sadamori Kojaku & & &  & \textbf{Time:} & TuTh 9:45 -- 11:15 \\
\textbf{Email:} &  \href{mailto:skojaku@binghamton.edu}{skojaku@binghamton.edu} & & & & \textbf{Place:} & J01 Engineering Bldg. \\
\textbf{Office:} &  J19 Engineering Bldg. & & & & \textbf{Credit:} & 3 \\
\textbf{Office Hour:} &  13:00 - 15:00 TuTh & & & & \textbf{Zoom:} & \href{https://binghamton.zoom.us/my/skojaku.zoom}{my/skojaku.zoom}
\end{tabular}
\end{minipage}
\rule{6in}{0.4pt}
\end{center}
\vspace{.5cm}
\setlength{\unitlength}{1in}
\renewcommand{\arraystretch}{2}

\noindent\textbf{Course Pages:} \begin{enumerate}
\item Course Website \url{https://skojaku.github.io/adv-net-sci}
\item Course Github \url{http://github.com/skojaku/adv-net-sci}
\end{enumerate}

\vskip.15in
\noindent\textbf{Office Hours:} Tuesday and Thursday 13:00 - 15:00.

\vskip.15in
\noindent\textbf{Main References:} %\footnotemark
Below is a curated list of essential books that will be referenced during the course.
\begin{itemize}
\item Mark Newman. {\textit{Networks (Second Edition)}}. Oxford University Press, 2018.
\item Filippo Menczer, Santo Fortunato, and Clayton A. Davis. {\textit{A First Course in Network Science}}. Cambridge University Press, 2020.
\item James Bagrow and Yong-Yeol Ahn. {\textit{Working with Network Data: A Data Science Perspective}}. Cambridge University Press, 2024.
\end{itemize}

% \footnotetext{Downloadable ebook versions are available on AeLP.}

\vskip.15in
\noindent\textbf{Objectives:}
Networks are all around us, from the vast expanse of the Internet to the intricate web of social connections that we build in our daily lives. But networks aren't just limited to the human realm---they can be found in every corner of the natural world, such as the complex interactions between animals, proteins, viruses, and DNAs. In recent years, we have witnessed remarkable advances in Al/ML and an ever-increasing volume and quality of data. Together, they offer an unprecedented opportunity to unlock the secrets of the world around us. This course is an introduction to network data analysis from the bottom up: through the interactions with network data and tools, we will learn how to store, manipulate, compute, and leverage network data in practice, as well as their underlying theoretical foundations.


\vskip.15in
\noindent\textbf{Expected Student Learning Outcomes:}
After completing this course, students will:
\begin{itemize}
    \item be able to interpret and evaluate modern network science literature, concepts, methodologies, tools, and recent research topics,
    \item be able to conduct advanced network modeling, analysis and simulation using appropriate mathematical/computational means,
    \item be able to design and conduct original research using network science methods and tools, and
    \item be able to demonstrate integration of Systems Science and/or Industrial and Systems Engineering knowledge and techniques and advanced network modeling/analysis/simulation skills in the form of a final project.
\end{itemize}

\vskip.15in
\noindent\textbf{Prerequisites:}
SSIE-523. Fluency in Python. Basic understanding of mathematics and statistics.

\clearpage

%\vspace*{.25in}
\noindent \textbf{Course Outline:} \textit{(MP = Mini Project: a hands-on group session)}
\begin{center}
\begin{minipage}{5in}
\begin{flushleft}
\textbf{August 2026:} \\
08/18 (Tue) \dotfill Seven Bridges of K\"onigsberg \\
08/20 (Thu) \dotfill Seven Bridges of K\"onigsberg \\
08/25 (Tue) \dotfill Small-world Networks \\
08/27 (Thu) \dotfill Small-world Networks \& Project Matchmaking\\

\textbf{September 2026:} \\
09/01 (Tue) \dotfill Small-world Networks \\
09/03 (Thu) \dotfill Small-world Networks / MP \\
09/08 (Tue) \dotfill No class (Monday classes meet) \\
09/10 (Thu) \dotfill Network Robustness \\
09/15 (Tue) \dotfill Network Robustness \\
09/17 (Thu) \dotfill Network Robustness / MP \\
09/22 (Tue) \dotfill Friendship Paradox \\
09/24 (Thu) \dotfill Friendship Paradox \\
09/29 (Tue) \dotfill Friendship Paradox / MP \\

\textbf{October 2026:} \\
10/01 (Thu) \dotfill Clustering \\
10/06 (Tue) \dotfill Clustering \\
10/08 (Thu) \dotfill Clustering / MP \\
10/13 (Tue) \dotfill No class (Fall break) \\
10/15 (Thu) \dotfill No class (Fall break) \\
10/20 (Tue) \dotfill Centrality \\
10/22 (Thu) \dotfill Centrality \\
10/27 (Tue) \dotfill Centrality / MP \\
10/29 (Thu) \dotfill Random Walks \\

\textbf{November 2026:} \\
11/03 (Tue) \dotfill Random Walks \\
11/05 (Thu) \dotfill Random Walks / MP \\
11/10 (Tue) \dotfill Graph Embedding \\
11/12 (Thu) \dotfill Graph Embedding \\
11/17 (Tue) \dotfill Graph Embedding / MP \\
11/19 (Thu) \dotfill Graph Neural Networks \\
11/24 (Tue) \dotfill Graph Neural Networks \\
11/26 (Thu) \dotfill No class (Thanksgiving) \\

\textbf{December 2026:} \\
12/01 (Tue) \dotfill Graph Neural Networks / MP \\
12/03 (Thu) \dotfill Presentation \\
12/08 (Tue) \dotfill Presentation \& wrap-up \\
Final Examinations \dotfill ~ 12/10-12/16 \\
\end{flushleft}
\end{minipage}
\end{center}

\vspace*{.15in}
\noindent\textbf{Grading Policy:}
\begin{itemize}
\item \textit{Grading Item}: Quiz (10\%), Student Lecture (10\%), Pair Notebook (10\%), Group Mini-Project (10\%), Exam (30\%), Project (30\%).
\item \textit{Bonus points}: None. The 100 points above are the whole grade.
\item \textit{Credits}: 3 credits.
\item \textit{Grading}: Normal grading; A through F.
\end{itemize}

\vskip.15in
\noindent\textbf{Course structure:}
``Don't think! Feeeeeel'' is a famous quote by Bruce Lee in the movie \textit{Enter the Dragon}, and this is the guiding philosophy of this course.
The primary goal of this course is to \textbf{feel} the concepts and tools of network science through pen-and-paper exercises and hands-on coding.
The course will first cover the practical skills of network science, followed by their theoretical foundations.
The course activities include:
\begin{itemize}
\item The class will begin with a weekly quiz on Brightspace about the lecture topics in the previous week. This quiz is intended to review the lecture topics and to test your understanding of the concepts and tools of network science.
\item Biweekly assignment on Github Classroom. Most are coding assignments, but not all of them are. The assignments are intended for practicing and equipping you with the skills to analyze network data.
\item A short lecture to cover the theoretical aspects.
\end{itemize}

\vskip.15in
\noindent\textbf{Communications:}
We use Discord for quicker informal communications, Q\&A, team discussions, and other casual conversations. We will send you an invitation link through Brightspace. Feel free to NOT use your full name (e.g., “Jane D.”)
Announcements will be sent via Brightspace and Discord. Many course-related information will be shared on Discord. So, you will miss a lot of information if you are not on Discord.

\vskip.15in
\noindent\textbf{Office Hours:}
\begin{itemize}
    \item \textbf{Sadamori Kojaku:} 13:00-15:00 Tuesdays and Thursdays at J19 Engineering Building (in person) or Zoom (\url{https://binghamton.zoom.us/my/skojaku.zoom}).
\end{itemize}

\clearpage
\noindent\textbf{Important Dates:}
\begin{center} \begin{minipage}{3.8in}
\begin{flushleft}
Project Proposal \#1      \dotfill ~09/27 (Sun) \\
Project Presentation \#2      \dotfill ~12/03 and 12/08 \\
Project Final Paper \#3      \dotfill ~12/06 (Sun) \\
Final Exam       \dotfill ~12/10-12/16 \\
\end{flushleft}
\end{minipage}
\end{center}

\vskip.15in
\noindent\textbf{Course Policy:}

\begin{itemize}
    \item \textbf{Attendance}: Residential students are expected to attend class in person; studies show that in-person attendance is associated with better learning outcomes [there]\footnote{\url{https://journals.sagepub.com/doi/10.3102/0034654310362998}}, [there]\footnote{\url{https://www.nature.com/articles/s41599-023-02590-1}}, [there]\footnote{\url{https://www.educationnext.org/zooming-to-class-slows-student-learning/}}). If you cannot attend, submit the absence request form (\url{https://forms.gle/yhzHoMSaKCRJXdGm7}) and email the instructor, preferably one day before class. \textbf{An absence not submitted through the form is not counted}, even if you emailed about it. Excused absences are limited to illness, accidents, job interviews, and conference travel. If you will miss more than two weeks for a legitimate reason, request an accommodation by email.
    \item \textbf{Laptop and mobile}: You may use laptops and mobile phones in class.
    \item \textbf{Credit hours}: This is a 3-credit course. In addition to the scheduled lectures, students are expected to spend at least 6.5 hours per week on course-related work, such as readings, lab sessions, studying for tests and examinations, and preparing assignments.
    \item \textbf{Generative AI}: You may use AI tools as learning aids for understanding course materials. However, the final submitted assignment must be your own work. If parts of an assignment are produced by generative AI, indicate the generated parts and cite the tools used, following this guideline: \url{https://style.mla.org/citing-generative-ai/}

    Try each assignment on your own first before turning to AI. Most assignments in this course can be solved by AI in seconds, but the purpose of an assignment is to build a skill, and that skill is only built by working through the problem yourself. If you understand the underlying concept, you can judge whether the AI's output is correct and use it effectively; if you do not, you will not be able to check its answers, and the same kind of problem will stop you again later.

    \item \textbf{Generative AI Service}: The following services are available to you. Google Gemini is free for all students; log in with your BU Google account at \url{https://gemini.google.com}. chat.binghamton.edu is a campus-hosted LLM chat service run by ITS; you can issue your own API key from the account settings and use it with coding agents (e.g., Claude Code, Cursor, Codex) at no cost. OpenRouter (\url{https://openrouter.ai}) offers a free tier of open-weight LLMs (50 requests/day); adding a one-time \$10 credit permanently raises this to 1{,}000 requests/day.
    \item \textbf{Data backup}: Back up your data and code regularly, for example with Google Drive, Dropbox, or GitHub (\url{https://github.com}). Loss of data or code (e.g., due to a laptop failure) is not an acceptable excuse for a delayed or missing submission.
    \item \textbf{Graduate Academic Consultants}: Binghamton University's Graduate Academic Consultants can help you with projects in this course. This is a free service. \url{https://www.binghamton.edu/grad-school/academic-support/graduate-academic-consultants/}
    \item \textbf{Disabilities}: Every attempt will be made to accommodate qualified students with disabilities (e.g. mental health, learning, chronic health, physical, hearing, vision, neurological, etc.). You must have established your eligibility for support services through Services for Students with Disabilities. Note that services are confidential, may take time to put into place, and are not retroactive. The office is located in the University Union, room 119. Captions and alternate media for print materials may take three or more weeks to get produced. Please contact Disability Services for Students at \url{https://www.binghamton.edu/ssd/index.html} or 607-777-2686 as soon as possible if accommodations are needed.
    \item \textbf{Bias-based incidents}: Any act of discrimination or harassment based on race, ethnicity, religious affiliation, gender, gender identity, sexual orientation, or disability can be reported at \url{https://www.binghamton.edu/diversity-equity-inclusion/reportbias.html} or to the Binghamton University Affirmative Action Officer at 607-777-4775.
    Sexual misconduct and Title IX
    \item \textbf{Sexual misconduct and Title IX}: Title IX and BU’s Sexual Harassment Policy regard any form of sexual harassment as a violation of the standards of conduct required of all persons associated with the institution. If you have experienced sexual misconduct or know someone who has, you can ask support from the University Counseling Center at 607-777-2772 (counseling, advocacy, and advice services). It is also important that you know that Title IX and University policy require me to share any information brought to my attention about potential sexual misconduct with the campus Deputy Title IX Coordinator or BU’s Title IX Coordinator. In that event, those individuals will work to ensure that appropriate measures are taken and resources are made available. Protecting student privacy is of utmost concern, and information will only be shared with those that need to know to ensure the University can respond and assist. Visit \url{https://www.binghamton.edu/counseling/resources/faculty/assault.html} and \url{https://www.binghamton.edu/services/title-ix/index.html} to learn more.
    \item \textbf{Mental health issues}: If you have any mental health issues, don’t hesitate to contact BU’s \href{https://www.binghamton.edu/counseling/index.html}{University Counseling Center}, which provides free counseling sessions. Also, please contact Disability Services for Students at \href{https://www.binghamton.edu/ssd/index.html}{Services for Students with Disabilities} or 607-777-6893 as soon as possible if accommodations are needed.
    \item \textbf{Academic Integrity}: Academic integrity is fundamental to the mission of our university. All students are expected to uphold the highest standards of academic honesty in their coursework and research. Violations of academic integrity include plagiarism, cheating, unauthorized collaboration, fabrication, and misrepresentation. For full details on the Academic Honesty Code, procedures, and appeals process, please refer to \href{https://www.binghamton.edu/watson/about/academic-honesty.html}{the Student Academic Honesty Code}. If you are unsure about what constitutes academic dishonesty in any situation, ask your instructor for clarification.
\end{itemize}

%%%%%% THE END
\end{document}
